% Options for packages loaded elsewhere
% Options for packages loaded elsewhere
\PassOptionsToPackage{unicode}{hyperref}
\PassOptionsToPackage{hyphens}{url}
\PassOptionsToPackage{dvipsnames,svgnames,x11names}{xcolor}
%
\documentclass[
  spanish,
  letterpaper,
  DIV=11,
  numbers=noendperiod]{scrreprt}
\usepackage{xcolor}
\usepackage{amsmath,amssymb}
\setcounter{secnumdepth}{5}
\usepackage{iftex}
\ifPDFTeX
  \usepackage[T1]{fontenc}
  \usepackage[utf8]{inputenc}
  \usepackage{textcomp} % provide euro and other symbols
\else % if luatex or xetex
  \usepackage{unicode-math} % this also loads fontspec
  \defaultfontfeatures{Scale=MatchLowercase}
  \defaultfontfeatures[\rmfamily]{Ligatures=TeX,Scale=1}
\fi
\usepackage{lmodern}
\ifPDFTeX\else
  % xetex/luatex font selection
\fi
% Use upquote if available, for straight quotes in verbatim environments
\IfFileExists{upquote.sty}{\usepackage{upquote}}{}
\IfFileExists{microtype.sty}{% use microtype if available
  \usepackage[]{microtype}
  \UseMicrotypeSet[protrusion]{basicmath} % disable protrusion for tt fonts
}{}
\makeatletter
\@ifundefined{KOMAClassName}{% if non-KOMA class
  \IfFileExists{parskip.sty}{%
    \usepackage{parskip}
  }{% else
    \setlength{\parindent}{0pt}
    \setlength{\parskip}{6pt plus 2pt minus 1pt}}
}{% if KOMA class
  \KOMAoptions{parskip=half}}
\makeatother
% Make \paragraph and \subparagraph free-standing
\makeatletter
\ifx\paragraph\undefined\else
  \let\oldparagraph\paragraph
  \renewcommand{\paragraph}{
    \@ifstar
      \xxxParagraphStar
      \xxxParagraphNoStar
  }
  \newcommand{\xxxParagraphStar}[1]{\oldparagraph*{#1}\mbox{}}
  \newcommand{\xxxParagraphNoStar}[1]{\oldparagraph{#1}\mbox{}}
\fi
\ifx\subparagraph\undefined\else
  \let\oldsubparagraph\subparagraph
  \renewcommand{\subparagraph}{
    \@ifstar
      \xxxSubParagraphStar
      \xxxSubParagraphNoStar
  }
  \newcommand{\xxxSubParagraphStar}[1]{\oldsubparagraph*{#1}\mbox{}}
  \newcommand{\xxxSubParagraphNoStar}[1]{\oldsubparagraph{#1}\mbox{}}
\fi
\makeatother

\usepackage{color}
\usepackage{fancyvrb}
\newcommand{\VerbBar}{|}
\newcommand{\VERB}{\Verb[commandchars=\\\{\}]}
\DefineVerbatimEnvironment{Highlighting}{Verbatim}{commandchars=\\\{\}}
% Add ',fontsize=\small' for more characters per line
\usepackage{framed}
\definecolor{shadecolor}{RGB}{241,243,245}
\newenvironment{Shaded}{\begin{snugshade}}{\end{snugshade}}
\newcommand{\AlertTok}[1]{\textcolor[rgb]{0.68,0.00,0.00}{#1}}
\newcommand{\AnnotationTok}[1]{\textcolor[rgb]{0.37,0.37,0.37}{#1}}
\newcommand{\AttributeTok}[1]{\textcolor[rgb]{0.40,0.45,0.13}{#1}}
\newcommand{\BaseNTok}[1]{\textcolor[rgb]{0.68,0.00,0.00}{#1}}
\newcommand{\BuiltInTok}[1]{\textcolor[rgb]{0.00,0.23,0.31}{#1}}
\newcommand{\CharTok}[1]{\textcolor[rgb]{0.13,0.47,0.30}{#1}}
\newcommand{\CommentTok}[1]{\textcolor[rgb]{0.37,0.37,0.37}{#1}}
\newcommand{\CommentVarTok}[1]{\textcolor[rgb]{0.37,0.37,0.37}{\textit{#1}}}
\newcommand{\ConstantTok}[1]{\textcolor[rgb]{0.56,0.35,0.01}{#1}}
\newcommand{\ControlFlowTok}[1]{\textcolor[rgb]{0.00,0.23,0.31}{\textbf{#1}}}
\newcommand{\DataTypeTok}[1]{\textcolor[rgb]{0.68,0.00,0.00}{#1}}
\newcommand{\DecValTok}[1]{\textcolor[rgb]{0.68,0.00,0.00}{#1}}
\newcommand{\DocumentationTok}[1]{\textcolor[rgb]{0.37,0.37,0.37}{\textit{#1}}}
\newcommand{\ErrorTok}[1]{\textcolor[rgb]{0.68,0.00,0.00}{#1}}
\newcommand{\ExtensionTok}[1]{\textcolor[rgb]{0.00,0.23,0.31}{#1}}
\newcommand{\FloatTok}[1]{\textcolor[rgb]{0.68,0.00,0.00}{#1}}
\newcommand{\FunctionTok}[1]{\textcolor[rgb]{0.28,0.35,0.67}{#1}}
\newcommand{\ImportTok}[1]{\textcolor[rgb]{0.00,0.46,0.62}{#1}}
\newcommand{\InformationTok}[1]{\textcolor[rgb]{0.37,0.37,0.37}{#1}}
\newcommand{\KeywordTok}[1]{\textcolor[rgb]{0.00,0.23,0.31}{\textbf{#1}}}
\newcommand{\NormalTok}[1]{\textcolor[rgb]{0.00,0.23,0.31}{#1}}
\newcommand{\OperatorTok}[1]{\textcolor[rgb]{0.37,0.37,0.37}{#1}}
\newcommand{\OtherTok}[1]{\textcolor[rgb]{0.00,0.23,0.31}{#1}}
\newcommand{\PreprocessorTok}[1]{\textcolor[rgb]{0.68,0.00,0.00}{#1}}
\newcommand{\RegionMarkerTok}[1]{\textcolor[rgb]{0.00,0.23,0.31}{#1}}
\newcommand{\SpecialCharTok}[1]{\textcolor[rgb]{0.37,0.37,0.37}{#1}}
\newcommand{\SpecialStringTok}[1]{\textcolor[rgb]{0.13,0.47,0.30}{#1}}
\newcommand{\StringTok}[1]{\textcolor[rgb]{0.13,0.47,0.30}{#1}}
\newcommand{\VariableTok}[1]{\textcolor[rgb]{0.07,0.07,0.07}{#1}}
\newcommand{\VerbatimStringTok}[1]{\textcolor[rgb]{0.13,0.47,0.30}{#1}}
\newcommand{\WarningTok}[1]{\textcolor[rgb]{0.37,0.37,0.37}{\textit{#1}}}

\usepackage{longtable,booktabs,array}
\newcounter{none} % for unnumbered tables
\usepackage{calc} % for calculating minipage widths
% Correct order of tables after \paragraph or \subparagraph
\usepackage{etoolbox}
\makeatletter
\patchcmd\longtable{\par}{\if@noskipsec\mbox{}\fi\par}{}{}
\makeatother
% Allow footnotes in longtable head/foot
\IfFileExists{footnotehyper.sty}{\usepackage{footnotehyper}}{\usepackage{footnote}}
\makesavenoteenv{longtable}
\usepackage{graphicx}
\makeatletter
\newsavebox\pandoc@box
\newcommand*\pandocbounded[1]{% scales image to fit in text height/width
  \sbox\pandoc@box{#1}%
  \Gscale@div\@tempa{\textheight}{\dimexpr\ht\pandoc@box+\dp\pandoc@box\relax}%
  \Gscale@div\@tempb{\linewidth}{\wd\pandoc@box}%
  \ifdim\@tempb\p@<\@tempa\p@\let\@tempa\@tempb\fi% select the smaller of both
  \ifdim\@tempa\p@<\p@\scalebox{\@tempa}{\usebox\pandoc@box}%
  \else\usebox{\pandoc@box}%
  \fi%
}
% Set default figure placement to htbp
\def\fps@figure{htbp}
\makeatother



\ifLuaTeX
\usepackage[bidi=basic,shorthands=off]{babel}
\else
\usepackage[bidi=default,shorthands=off]{babel}
\fi
\ifLuaTeX
  \usepackage{selnolig} % disable illegal ligatures
\fi


\setlength{\emergencystretch}{3em} % prevent overfull lines

\providecommand{\tightlist}{%
  \setlength{\itemsep}{0pt}\setlength{\parskip}{0pt}}



 


\usepackage{fvextra}
\RecustomVerbatimEnvironment{verbatim}{Verbatim}{
  frame=single, rulecolor=\color{black!40},
  bgcolor=black!8, fontsize=\small, framesep=3mm}
\KOMAoption{captions}{tableheading}
\makeatletter
\@ifpackageloaded{tcolorbox}{}{\usepackage[skins,breakable]{tcolorbox}}
\@ifpackageloaded{fontawesome5}{}{\usepackage{fontawesome5}}
\definecolor{quarto-callout-color}{HTML}{909090}
\definecolor{quarto-callout-note-color}{HTML}{0758E5}
\definecolor{quarto-callout-important-color}{HTML}{CC1914}
\definecolor{quarto-callout-warning-color}{HTML}{EB9113}
\definecolor{quarto-callout-tip-color}{HTML}{00A047}
\definecolor{quarto-callout-caution-color}{HTML}{FC5300}
\definecolor{quarto-callout-color-frame}{HTML}{acacac}
\definecolor{quarto-callout-note-color-frame}{HTML}{4582ec}
\definecolor{quarto-callout-important-color-frame}{HTML}{d9534f}
\definecolor{quarto-callout-warning-color-frame}{HTML}{f0ad4e}
\definecolor{quarto-callout-tip-color-frame}{HTML}{02b875}
\definecolor{quarto-callout-caution-color-frame}{HTML}{fd7e14}
\makeatother
\makeatletter
\@ifpackageloaded{bookmark}{}{\usepackage{bookmark}}
\makeatother
\makeatletter
\@ifpackageloaded{caption}{}{\usepackage{caption}}
\AtBeginDocument{%
\ifdefined\contentsname
  \renewcommand*\contentsname{Tabla de contenidos}
\else
  \newcommand\contentsname{Tabla de contenidos}
\fi
\ifdefined\listfigurename
  \renewcommand*\listfigurename{Listado de Figuras}
\else
  \newcommand\listfigurename{Listado de Figuras}
\fi
\ifdefined\listtablename
  \renewcommand*\listtablename{Listado de Tablas}
\else
  \newcommand\listtablename{Listado de Tablas}
\fi
\ifdefined\figurename
  \renewcommand*\figurename{Figura}
\else
  \newcommand\figurename{Figura}
\fi
\ifdefined\tablename
  \renewcommand*\tablename{Tabla}
\else
  \newcommand\tablename{Tabla}
\fi
}
\@ifpackageloaded{float}{}{\usepackage{float}}
\floatstyle{ruled}
\@ifundefined{c@chapter}{\newfloat{codelisting}{h}{lop}}{\newfloat{codelisting}{h}{lop}[chapter]}
\floatname{codelisting}{Listado}
\newcommand*\listoflistings{\listof{codelisting}{Listado de Listados}}
\makeatother
\makeatletter
\makeatother
\makeatletter
\@ifpackageloaded{caption}{}{\usepackage{caption}}
\@ifpackageloaded{subcaption}{}{\usepackage{subcaption}}
\makeatother
\usepackage{bookmark}
\IfFileExists{xurl.sty}{\usepackage{xurl}}{} % add URL line breaks if available
\urlstyle{same}
% fallback for those not using the hyperref driver hyperxmp:
\makeatletter
\@ifundefined{xmpquote}{\newcommand{\xmpquote}[1]{#1}}{}
\makeatother
\hypersetup{
  pdftitle={Introducción a Python para Ingeniería{,} Datos y Economía},
  pdfauthor={John Leal},
  pdflang={es},
  colorlinks=true,
  linkcolor={blue},
  filecolor={Maroon},
  citecolor={Blue},
  urlcolor={Blue},
  pdfcreator={LaTeX via pandoc}}


\title{Introducción a Python para Ingeniería, Datos y Economía}
\author{John Leal}
\date{2026-09-22}
\begin{document}
\maketitle

\renewcommand*\contentsname{Tabla de contenidos}
{
\setcounter{tocdepth}{2}
\tableofcontents
}

\bookmarksetup{startatroot}

\chapter*{Prólogo}\label{pruxf3logo}
\addcontentsline{toc}{chapter}{Prólogo}

\markboth{Prólogo}{Prólogo}

Este libro está pensado para estudiantes y profesionales de ingeniería,
análisis de datos y economía que desean aprender a programar en Python
desde cero. No se asume ninguna experiencia previa en programación: los
conceptos se introducen de forma progresiva, con ejemplos cercanos a la
ingeniería, la estadística y la economía, para que cada tema tenga un
contexto claro y aplicable.

Todo el código está diseñado para ejecutarse en Google Colab, el entorno
gratuito de notebooks de Google. No necesitas instalar nada en tu
computadora: basta un navegador y una cuenta de Google. Cada capítulo
contiene bloques de código listos para copiar y pegar en una celda de
Colab, de modo que puedes leer, ejecutar y experimentar de inmediato.

Al finalizar el libro serás capaz de escribir programas en Python,
procesar datos con pandas, resolver problemas de ingeniería y economía
(valor futuro, VAN, simulación Monte Carlo) y entender los conceptos
básicos del mundo del big data.

\section*{Cómo usar este libro}\label{cuxf3mo-usar-este-libro}
\addcontentsline{toc}{section}{Cómo usar este libro}

\markright{Cómo usar este libro}

\begin{enumerate}
\def\labelenumi{\arabic{enumi}.}
\tightlist
\item
  \textbf{Lee con Colab abierto.} Cada bloque de código está hecho para
  copiarse y pegarse directamente en una celda de Colab. Ejecuta todo lo
  que veas: programar se aprende haciendo.
\item
  \textbf{Modifica los ejemplos.} Cambia los valores, rompe el código a
  propósito, lee los errores. Experimentar es parte del método.
\item
  \textbf{Haz los ejercicios.} Al final de cada capítulo hay ejercicios
  y un mini-reto que integra lo aprendido con un problema de ingeniería,
  administración o economía.
\item
  \textbf{Lleva un notebook por capítulo.} Serán tu cuaderno de apuntes
  ejecutable y tu portafolio de práctica.
\end{enumerate}

La ruta del libro es progresiva: los capítulos 1 a 4 construyen los
fundamentos del lenguaje; el 5 y el 6 introducen el cómputo científico y
el manejo de datos; el 7 aplica todo a problemas de administración y
economía; y el 8 cierra con una introducción a big data.

\bookmarksetup{startatroot}

\chapter{Introducción a Python y Google
Colab}\label{introducciuxf3n-a-python-y-google-colab}

\section{¿Qué es Python y por qué lo usaremos en Google
Colab?}\label{quuxe9-es-python-y-por-quuxe9-lo-usaremos-en-google-colab}

Python es un lenguaje de programación de propósito general, creado con
una idea central: que el código se lea casi como texto en inglés. Esa
claridad lo convirtió en el lenguaje más usado del mundo en ciencia de
datos, ingeniería, economía e inteligencia artificial.

¿Por qué Python y no otro lenguaje? Porque con las mismas diez líneas de
código que en otros lenguajes apenas configurarían el entorno, en Python
ya estás analizando datos reales, graficando una función o calculando el
valor presente neto de un proyecto.

En este libro trabajaremos con \textbf{Google Colab}: un entorno
gratuito de Google que corre en el navegador. No necesitas instalar nada
en tu computador; solo una cuenta de Google. Cada bloque de código de
este libro está diseñado para que lo copies, lo pegues en una celda de
Colab y lo ejecutes.

\section{Tu primer programa}\label{tu-primer-programa}

\begin{enumerate}
\def\labelenumi{\arabic{enumi}.}
\tightlist
\item
  Abre
  \href{https://colab.research.google.com}{colab.research.google.com} en
  tu navegador.
\item
  Haz clic en \textbf{Nuevo notebook}.
\item
  En la primera celda, escribe (o copia) lo siguiente:
\end{enumerate}

\begin{Shaded}
\begin{Highlighting}[]
\BuiltInTok{print}\NormalTok{(}\StringTok{"Hola, mundo"}\NormalTok{)}
\end{Highlighting}
\end{Shaded}

\begin{verbatim}
Hola, mundo
\end{verbatim}

\begin{enumerate}
\def\labelenumi{\arabic{enumi}.}
\setcounter{enumi}{3}
\tightlist
\item
  Ejecuta la celda con el botón ▶️ o con las teclas
  \texttt{Shift\ +\ Enter}.
\end{enumerate}

¿Qué acaba de pasar? \texttt{print()} es una \textbf{función}: una
instrucción que le dice a Python ``muestra esto en pantalla''. El texto
entre comillas se llama \textbf{cadena de caracteres} (\emph{string}).
Acabas de escribir y ejecutar tu primer programa.

\begin{tcolorbox}[enhanced jigsaw, arc=.35mm, bottomrule=.15mm, bottomtitle=1mm, breakable, colback=white, colbacktitle=quarto-callout-note-color!10!white, colframe=quarto-callout-note-color-frame, coltitle=black, left=2mm, leftrule=.75mm, opacityback=0, opacitybacktitle=0.6, rightrule=.15mm, title=\textcolor{quarto-callout-note-color}{\faInfo}\hspace{0.5em}{Nota}, titlerule=0mm, toprule=.15mm, toptitle=1mm]

Un notebook de Colab tiene dos tipos de celdas: \textbf{código} (se
ejecuta) y \textbf{texto} (se lee, como este párrafo). Agrégalas con los
botones \texttt{+\ Código} y \texttt{+\ Texto}.

\end{tcolorbox}

\section{Python como calculadora}\label{python-como-calculadora}

Antes de programar ``en serio'', usa Python como una calculadora. Copia
y ejecuta cada línea:

\begin{Shaded}
\begin{Highlighting}[]
\DecValTok{2} \OperatorTok{+} \DecValTok{3}
\end{Highlighting}
\end{Shaded}

\begin{Shaded}
\begin{Highlighting}[]
\DecValTok{10} \OperatorTok{/} \DecValTok{4}
\end{Highlighting}
\end{Shaded}

\begin{Shaded}
\begin{Highlighting}[]
\DecValTok{2} \OperatorTok{**} \DecValTok{10}
\end{Highlighting}
\end{Shaded}

\begin{Shaded}
\begin{Highlighting}[]
\DecValTok{17} \OperatorTok{\%} \DecValTok{5}
\end{Highlighting}
\end{Shaded}

Los operadores básicos son:

\begin{longtable}[]{@{}clcc@{}}
\caption{Operadores aritméticos básicos de
Python}\label{tbl-operadores}\tabularnewline
\toprule\noalign{}
Operador & Significado & Ejemplo & Resultado \\
\midrule\noalign{}
\endfirsthead
\toprule\noalign{}
Operador & Significado & Ejemplo & Resultado \\
\midrule\noalign{}
\endhead
\bottomrule\noalign{}
\endlastfoot
\texttt{+} & Suma & \texttt{2\ +\ 3} & \texttt{5} \\
\texttt{-} & Resta & \texttt{7\ -\ 4} & \texttt{3} \\
\texttt{*} & Multiplicación & \texttt{3\ *\ 5} & \texttt{15} \\
\texttt{/} & División & \texttt{10\ /\ 4} & \texttt{2.5} \\
\texttt{**} & Potencia & \texttt{2\ **\ 10} & \texttt{1024} \\
\texttt{//} & División entera & \texttt{17\ //\ 5} & \texttt{3} \\
\texttt{\%} & Residuo (módulo) & \texttt{17\ \%\ 5} & \texttt{2} \\
\end{longtable}

Un ejemplo con sentido económico: si inviertes \$1.000.000 a una tasa
del 12 \% anual durante 5 años, el valor futuro es:

\begin{Shaded}
\begin{Highlighting}[]
\DecValTok{1000000} \OperatorTok{*}\NormalTok{ (}\DecValTok{1} \OperatorTok{+} \FloatTok{0.12}\NormalTok{) }\OperatorTok{**} \DecValTok{5}
\end{Highlighting}
\end{Shaded}

Observa que Python respeta el orden de las operaciones (primero
paréntesis, luego potencias, luego multiplicación y división, al final
suma y resta), igual que en álgebra.

\section{Los errores son tus amigos}\label{los-errores-son-tus-amigos}

Escribe esto en una celda y ejecútalo:

\begin{Shaded}
\begin{Highlighting}[]
\BuiltInTok{print}\NormalTok{(}\StringTok{"Hola"}
\end{Highlighting}
\end{Shaded}

Verás algo como:

\begin{Shaded}
\begin{Highlighting}[]
\NormalTok{SyntaxError: incomplete input}
\end{Highlighting}
\end{Shaded}

No te alarmes: un error no rompe nada. Python te está diciendo
\textbf{qué} falló (\texttt{SyntaxError}: error de sintaxis) y en este
caso es evidente: falta cerrar el paréntesis. Corrige y ejecuta de
nuevo.

\begin{tcolorbox}[enhanced jigsaw, arc=.35mm, bottomrule=.15mm, bottomtitle=1mm, breakable, colback=white, colbacktitle=quarto-callout-tip-color!10!white, colframe=quarto-callout-tip-color-frame, coltitle=black, left=2mm, leftrule=.75mm, opacityback=0, opacitybacktitle=0.6, rightrule=.15mm, title=\textcolor{quarto-callout-tip-color}{\faLightbulb}\hspace{0.5em}{Tip}, titlerule=0mm, toprule=.15mm, toptitle=1mm]

Regla de oro para todo el libro: \textbf{lee los errores de abajo hacia
arriba}. La última línea del mensaje casi siempre dice exactamente qué
salió mal.

\end{tcolorbox}

\section{Guardar tu trabajo}\label{guardar-tu-trabajo}

Colab guarda tus notebooks automáticamente en tu Google Drive (carpeta
\texttt{Colab\ Notebooks}). Buenas prácticas desde el primer día:

\begin{itemize}
\tightlist
\item
  Renombra el notebook: clic en el título de arriba a la izquierda. Usa
  nombres sin espacios ni tildes:
  \texttt{capitulo01\_introduccion.ipynb}.
\item
  Para descargarlo:
  \texttt{Archivo\ →\ Descargar\ →\ Descargar\ .ipynb}.
\item
  Crea un notebook por capítulo de este libro; te servirán como tu
  propio cuaderno de apuntes ejecutable.
\end{itemize}

\section{Ejercicios}\label{ejercicios}

\begin{enumerate}
\def\labelenumi{\arabic{enumi}.}
\tightlist
\item
  Ejecuta \texttt{print("Tu\ nombre")} con tu nombre real.
\item
  Calcula cuántos segundos tiene un año (365 días) usando solo
  operadores.
\item
  Calcula el valor futuro de \$5.000.000 al 8 \% anual durante 10 años.
\item
  Provoca a propósito un error quitando una comilla en un
  \texttt{print}, lee el mensaje y corrígelo.
\item
  \textbf{Mini-reto:} una hectárea produce 4,2 toneladas de café.
  ¿Cuántas toneladas producen 37 hectáreas? ¿Y si el precio es
  \$9.500.000 por tonelada, cuánto vale la cosecha total? Resuélvelo en
  una sola celda con dos operaciones.
\end{enumerate}

\section{Resumen}\label{resumen}

\begin{itemize}
\tightlist
\item
  Python es un lenguaje claro y es el estándar en datos e ingeniería.
\item
  Colab te permite programar sin instalar nada, desde el navegador.
\item
  \texttt{print()} muestra resultados en pantalla.
\item
  Python funciona como calculadora con los operadores
  \texttt{+\ -\ *\ /\ **\ //\ \%}.
\item
  Los errores son mensajes de ayuda: léelos de abajo hacia arriba.
\end{itemize}

\bookmarksetup{startatroot}

\chapter{Fundamentos: variables y tipos de
datos}\label{fundamentos-variables-y-tipos-de-datos}

\section{¿Qué es una variable?}\label{quuxe9-es-una-variable}

Imagina una caja con una etiqueta. Dentro de la caja guardas algo (un
número, un texto), y la etiqueta te permite encontrarla después. Eso es
una \textbf{variable}: un nombre que guarda un valor.

En Python, guardar un valor en una variable se llama \textbf{asignar}, y
se hace con el signo \texttt{=}:

\begin{Shaded}
\begin{Highlighting}[]
\NormalTok{precio }\OperatorTok{=} \DecValTok{9500}
\end{Highlighting}
\end{Shaded}

Lee esa línea así: ``la variable \texttt{precio} guarda el valor 9500''.
Ojo: el \texttt{=} de Python \textbf{no es} el igual de las matemáticas;
no afirma que dos cosas sean iguales, sino que \textbf{guarda} el valor
de la derecha en el nombre de la izquierda.

A partir de ahí, puedes usar la variable en cualquier cálculo:

\begin{Shaded}
\begin{Highlighting}[]
\NormalTok{precio }\OperatorTok{=} \DecValTok{9500}
\NormalTok{cantidad }\OperatorTok{=} \DecValTok{37}
\NormalTok{total }\OperatorTok{=}\NormalTok{ precio }\OperatorTok{*}\NormalTok{ cantidad}
\BuiltInTok{print}\NormalTok{(total)}
\end{Highlighting}
\end{Shaded}

La gran ventaja: si el precio cambia, solo modificas una línea y todo lo
demás se recalcula solo. Cambia \texttt{precio} a \texttt{10200} y
vuelve a ejecutar.

\section{Reglas para nombrar
variables}\label{reglas-para-nombrar-variables}

Python es estricto con los nombres. Estos son válidos:

\begin{Shaded}
\begin{Highlighting}[]
\NormalTok{precio\_tonelada }\OperatorTok{=} \DecValTok{9500}
\NormalTok{temperatura2 }\OperatorTok{=} \FloatTok{25.4}
\NormalTok{nombre\_finca }\OperatorTok{=} \StringTok{"La Esperanza"}
\end{Highlighting}
\end{Shaded}

Estos \textbf{no} son válidos (ejecútalos y lee el error):

\begin{Shaded}
\begin{Highlighting}[]
\DecValTok{2}\ErrorTok{temperatura} \OperatorTok{=} \DecValTok{25}      \CommentTok{\# no puede empezar con número}
\NormalTok{precio}\OperatorTok{{-}}\NormalTok{tonelada }\OperatorTok{=} \DecValTok{9500} \CommentTok{\# no puede tener guiones}
\NormalTok{precio tonelada }\OperatorTok{=} \DecValTok{9500} \CommentTok{\# no puede tener espacios}
\end{Highlighting}
\end{Shaded}

\begin{tcolorbox}[enhanced jigsaw, arc=.35mm, bottomrule=.15mm, bottomtitle=1mm, breakable, colback=white, colbacktitle=quarto-callout-note-color!10!white, colframe=quarto-callout-note-color-frame, coltitle=black, left=2mm, leftrule=.75mm, opacityback=0, opacitybacktitle=0.6, rightrule=.15mm, title=\textcolor{quarto-callout-note-color}{\faInfo}\hspace{0.5em}{Nota}, titlerule=0mm, toprule=.15mm, toptitle=1mm]

\textbf{Convención de este libro:} nombres en minúsculas, sin tildes,
con palabras separadas por guion bajo: \texttt{valor\_futuro},
\texttt{tasa\_anual}, \texttt{numero\_hectareas}. No es obligatorio,
pero es el estilo que usa todo el mundo de Python.

\end{tcolorbox}

\section{Los cuatro tipos de datos
básicos}\label{los-cuatro-tipos-de-datos-buxe1sicos}

Todo valor en Python tiene un \textbf{tipo}. Para empezar solo necesitas
conocer cuatro:

{\def\LTcaptype{none} % do not increment counter
\begin{longtable}[]{@{}clll@{}}
\toprule\noalign{}
Tipo & Nombre & Qué guarda & Ejemplo \\
\midrule\noalign{}
\endhead
\bottomrule\noalign{}
\endlastfoot
\texttt{int} & Entero & Números sin decimales & \texttt{37},
\texttt{-5}, \texttt{1000000} \\
\texttt{float} & Decimal & Números con decimales & \texttt{4.2},
\texttt{0.12}, \texttt{-3.5} \\
\texttt{str} & Texto & Cadenas de caracteres & \texttt{"café"},
\texttt{"La\ Esperanza"} \\
\texttt{bool} & Booleano & Verdadero o falso & \texttt{True},
\texttt{False} \\
\end{longtable}
}

Python detecta el tipo automáticamente según lo que guardes. Puedes
preguntárselo con la función \texttt{type()}:

\begin{Shaded}
\begin{Highlighting}[]
\NormalTok{precio }\OperatorTok{=} \DecValTok{9500}
\BuiltInTok{print}\NormalTok{(}\BuiltInTok{type}\NormalTok{(precio))}
\end{Highlighting}
\end{Shaded}

\begin{Shaded}
\begin{Highlighting}[]
\NormalTok{rendimiento }\OperatorTok{=} \FloatTok{4.2}
\BuiltInTok{print}\NormalTok{(}\BuiltInTok{type}\NormalTok{(rendimiento))}
\end{Highlighting}
\end{Shaded}

\begin{Shaded}
\begin{Highlighting}[]
\NormalTok{nombre }\OperatorTok{=} \StringTok{"Ana"}
\BuiltInTok{print}\NormalTok{(}\BuiltInTok{type}\NormalTok{(nombre))}
\end{Highlighting}
\end{Shaded}

\begin{Shaded}
\begin{Highlighting}[]
\NormalTok{es\_rentable }\OperatorTok{=} \VariableTok{True}
\BuiltInTok{print}\NormalTok{(}\BuiltInTok{type}\NormalTok{(es\_rentable))}
\end{Highlighting}
\end{Shaded}

\begin{tcolorbox}[enhanced jigsaw, arc=.35mm, bottomrule=.15mm, bottomtitle=1mm, breakable, colback=white, colbacktitle=quarto-callout-warning-color!10!white, colframe=quarto-callout-warning-color-frame, coltitle=black, left=2mm, leftrule=.75mm, opacityback=0, opacitybacktitle=0.6, rightrule=.15mm, title=\textcolor{quarto-callout-warning-color}{\faExclamationTriangle}\hspace{0.5em}{Advertencia}, titlerule=0mm, toprule=.15mm, toptitle=1mm]

\textbf{Cuidado con el punto y la coma:} en Python los decimales se
escriben con \textbf{punto}, no con coma. Escribe \texttt{4.2}, no
\texttt{4,2}. Si escribes \texttt{4,2} Python no dará un error obvio:
pensará que son dos números separados. Es uno de los errores más comunes
de quienes empiezan en países hispanos.

\end{tcolorbox}

\section{Convertir entre tipos}\label{convertir-entre-tipos}

A veces necesitas transformar un tipo en otro. Se hace con funciones que
tienen el nombre del tipo: \texttt{int()}, \texttt{float()},
\texttt{str()}.

El caso clásico: los números que vienen como texto no se pueden sumar:

\begin{Shaded}
\begin{Highlighting}[]
\NormalTok{a }\OperatorTok{=} \StringTok{"2"}
\NormalTok{b }\OperatorTok{=} \StringTok{"3"}
\BuiltInTok{print}\NormalTok{(a }\OperatorTok{+}\NormalTok{ b)}
\end{Highlighting}
\end{Shaded}

¿Sorprendido? El resultado es \texttt{"23"}, no \texttt{5}. Como
\texttt{a} y \texttt{b} son \textbf{texto}, el \texttt{+} los une en
lugar de sumarlos (pegar textos se llama \textbf{concatenar}). La
solución es convertirlos primero:

\begin{Shaded}
\begin{Highlighting}[]
\NormalTok{a }\OperatorTok{=} \StringTok{"2"}
\NormalTok{b }\OperatorTok{=} \StringTok{"3"}
\BuiltInTok{print}\NormalTok{(}\BuiltInTok{int}\NormalTok{(a) }\OperatorTok{+} \BuiltInTok{int}\NormalTok{(b))}
\end{Highlighting}
\end{Shaded}

\section{\texorpdfstring{Pedir datos al usuario:
\texttt{input()}}{Pedir datos al usuario: input()}}\label{pedir-datos-al-usuario-input}

Hasta ahora los valores estaban escritos en el código. Pero normalmente
querrás preguntarle a la persona que usa tu programa. Eso se hace con
\texttt{input()}:

\begin{Shaded}
\begin{Highlighting}[]
\NormalTok{nombre }\OperatorTok{=} \BuiltInTok{input}\NormalTok{(}\StringTok{"¿Cómo te llamas? "}\NormalTok{)}
\BuiltInTok{print}\NormalTok{(}\StringTok{"Hola,"}\NormalTok{, nombre)}
\end{Highlighting}
\end{Shaded}

Al ejecutar, el programa se detiene, muestra la pregunta y espera a que
escribas algo y presiones \texttt{Enter}.

\begin{tcolorbox}[enhanced jigsaw, arc=.35mm, bottomrule=.15mm, bottomtitle=1mm, breakable, colback=white, colbacktitle=quarto-callout-important-color!10!white, colframe=quarto-callout-important-color-frame, coltitle=black, left=2mm, leftrule=.75mm, opacityback=0, opacitybacktitle=0.6, rightrule=.15mm, title=\textcolor{quarto-callout-important-color}{\faExclamation}\hspace{0.5em}{Importante}, titlerule=0mm, toprule=.15mm, toptitle=1mm]

\textbf{Lo que input() devuelve siempre es texto} (\texttt{str}), aunque
la persona escriba un número. Si vas a hacer cuentas con esa respuesta,
conviértela:

\begin{Shaded}
\begin{Highlighting}[]
\NormalTok{edad }\OperatorTok{=} \BuiltInTok{input}\NormalTok{(}\StringTok{"¿Cuántos años tienes? "}\NormalTok{)}
\NormalTok{edad }\OperatorTok{=} \BuiltInTok{int}\NormalTok{(edad)}
\BuiltInTok{print}\NormalTok{(}\StringTok{"El próximo año tendrás"}\NormalTok{, edad }\OperatorTok{+} \DecValTok{1}\NormalTok{)}
\end{Highlighting}
\end{Shaded}

O en una sola línea, que es como lo verás en la práctica:

\begin{Shaded}
\begin{Highlighting}[]
\NormalTok{edad }\OperatorTok{=} \BuiltInTok{int}\NormalTok{(}\BuiltInTok{input}\NormalTok{(}\StringTok{"¿Cuántos años tienes? "}\NormalTok{))}
\end{Highlighting}
\end{Shaded}

\end{tcolorbox}

\section{Imprimir con formato: las
f-strings}\label{imprimir-con-formato-las-f-strings}

Ya sabes que \texttt{print()} muestra cosas. Cuando quieras mezclar
texto con variables de forma elegante, usa una \textbf{f-string}: una
cadena con una \texttt{f} antes de las comillas, donde las variables van
entre llaves \texttt{\{\}}:

\begin{Shaded}
\begin{Highlighting}[]
\NormalTok{precio }\OperatorTok{=} \DecValTok{9500}
\NormalTok{cantidad }\OperatorTok{=} \DecValTok{37}
\NormalTok{total }\OperatorTok{=}\NormalTok{ precio }\OperatorTok{*}\NormalTok{ cantidad}
\BuiltInTok{print}\NormalTok{(}\SpecialStringTok{f"El total de la cosecha es $}\SpecialCharTok{\{}\NormalTok{total}\SpecialCharTok{\}}\SpecialStringTok{"}\NormalTok{)}
\end{Highlighting}
\end{Shaded}

Todo lo que está dentro de las llaves se evalúa. Puedes incluso hacer
cuentas ahí dentro:

\begin{Shaded}
\begin{Highlighting}[]
\NormalTok{tasa }\OperatorTok{=} \FloatTok{0.12}
\NormalTok{capital }\OperatorTok{=} \DecValTok{1000000}
\BuiltInTok{print}\NormalTok{(}\SpecialStringTok{f"En 5 años tendrás $}\SpecialCharTok{\{}\NormalTok{capital }\OperatorTok{*}\NormalTok{ (}\DecValTok{1} \OperatorTok{+}\NormalTok{ tasa) }\OperatorTok{**} \DecValTok{5}\SpecialCharTok{\}}\SpecialStringTok{"}\NormalTok{)}
\end{Highlighting}
\end{Shaded}

\section{Comentarios: notas para
humanos}\label{comentarios-notas-para-humanos}

Todo lo que escribas después de un \texttt{\#} en una línea, Python lo
ignora por completo. Sirve para dejar notas que expliquen tu código:

\begin{Shaded}
\begin{Highlighting}[]
\CommentTok{\# Cálculo del valor futuro a 5 años}
\NormalTok{tasa }\OperatorTok{=} \FloatTok{0.12}          \CommentTok{\# tasa anual}
\NormalTok{capital }\OperatorTok{=} \DecValTok{1000000}    \CommentTok{\# inversión inicial en pesos}
\NormalTok{vf }\OperatorTok{=}\NormalTok{ capital }\OperatorTok{*}\NormalTok{ (}\DecValTok{1} \OperatorTok{+}\NormalTok{ tasa) }\OperatorTok{**} \DecValTok{5}
\BuiltInTok{print}\NormalTok{(}\SpecialStringTok{f"Valor futuro: $}\SpecialCharTok{\{}\NormalTok{vf}\SpecialCharTok{\}}\SpecialStringTok{"}\NormalTok{)}
\end{Highlighting}
\end{Shaded}

Comenta tu código pensando en la persona que lo leerá en seis meses. Esa
persona, casi siempre, serás tú.

\section{Ejercicios}\label{ejercicios-1}

\begin{enumerate}
\def\labelenumi{\arabic{enumi}.}
\item
  Crea tres variables con tu nombre, tu edad y tu estatura (en metros,
  con decimal). Imprímelas con una f-string en una sola frase.
\item
  Pregunta al usuario cuántas hectáreas tiene su finca y cuántas
  toneladas produce por hectárea. Calcula e imprime la producción total.
\item
  Pregunta al usuario su salario mensual y calcula cuánto gana al año.
  Imprime el resultado con una f-string.
\item
  Ejecuta esto, lee el error y corrígelo:

\begin{Shaded}
\begin{Highlighting}[]
\NormalTok{precio }\OperatorTok{=} \BuiltInTok{input}\NormalTok{(}\StringTok{"Precio: "}\NormalTok{)}
\BuiltInTok{print}\NormalTok{(precio }\OperatorTok{*} \DecValTok{2}\NormalTok{)}
\end{Highlighting}
\end{Shaded}

  (Pista: multiplica texto por dos y verás algo raro, pero no es lo que
  queremos).
\item
  \textbf{Mini-reto (economía):} pide al usuario un capital inicial, una
  tasa de interés anual (en porcentaje, ej. 12) y un número de años.
  Calcula el valor futuro con interés compuesto e imprímelo con formato:
  \texttt{Con\ \$X\ al\ Y\%\ durante\ Z\ años,\ tendrás\ \$W}. Recuerda
  convertir la tasa de porcentaje a decimal dividiendo entre 100.
\end{enumerate}

\section{Resumen}\label{resumen-1}

\begin{itemize}
\tightlist
\item
  Una \textbf{variable} es un nombre que guarda un valor; se asigna con
  \texttt{=}.
\item
  Los cuatro tipos básicos: \texttt{int} (enteros), \texttt{float}
  (decimales), \texttt{str} (texto), \texttt{bool} (verdadero/falso).
\item
  Los decimales se escriben con \textbf{punto}: \texttt{4.2}.
\item
  \texttt{input()} siempre devuelve texto; usa \texttt{int()} o
  \texttt{float()} para hacer cuentas.
\item
  Las \textbf{f-strings} (\texttt{f"...\{variable\}..."}) mezclan texto
  y valores fácilmente.
\item
  Los \textbf{comentarios} con \texttt{\#} documentan tu código.
\end{itemize}

\bookmarksetup{startatroot}

\chapter{Control de flujo: decisiones, repeticiones y
funciones}\label{control-de-flujo-decisiones-repeticiones-y-funciones}

Hasta ahora tus programas ejecutan las líneas una tras otra, de arriba
hacia abajo, siempre igual. Los programas útiles hacen tres cosas más:
\textbf{deciden} (si pasa esto, haz aquello), \textbf{repiten} (haz esto
1000 veces) y \textbf{empaquetan} (guarda este procedimiento para usarlo
cuando quiera). Eso es el control de flujo.

\section{\texorpdfstring{Decisiones: \texttt{if}, \texttt{elif},
\texttt{else}}{Decisiones: if, elif, else}}\label{decisiones-if-elif-else}

Un sistema de riego automático tiene una lógica muy simple: \emph{si la
humedad del suelo es baja, activar la bomba}. En Python eso se escribe
casi igual que en español:

\begin{Shaded}
\begin{Highlighting}[]
\NormalTok{humedad }\OperatorTok{=} \DecValTok{25}  \CommentTok{\# porcentaje medido por el sensor}

\ControlFlowTok{if}\NormalTok{ humedad }\OperatorTok{\textless{}} \DecValTok{30}\NormalTok{:}
    \BuiltInTok{print}\NormalTok{(}\StringTok{"Humedad baja: activar riego"}\NormalTok{)}
\end{Highlighting}
\end{Shaded}

Ejecútalo. Luego cambia \texttt{humedad} a \texttt{45} y ejecútalo de
nuevo: no aparece nada, porque la condición ya no se cumple.

Para cubrir más casos se usan \texttt{elif} (``si no, entonces
si\ldots{}'') y \texttt{else} (``en cualquier otro caso''):

\begin{Shaded}
\begin{Highlighting}[]
\NormalTok{humedad }\OperatorTok{=} \DecValTok{25}

\ControlFlowTok{if}\NormalTok{ humedad }\OperatorTok{\textless{}} \DecValTok{30}\NormalTok{:}
    \BuiltInTok{print}\NormalTok{(}\StringTok{"Humedad baja: activar riego"}\NormalTok{)}
\ControlFlowTok{elif}\NormalTok{ humedad }\OperatorTok{\textless{}} \DecValTok{60}\NormalTok{:}
    \BuiltInTok{print}\NormalTok{(}\StringTok{"Humedad adecuada: no hacer nada"}\NormalTok{)}
\ControlFlowTok{else}\NormalTok{:}
    \BuiltInTok{print}\NormalTok{(}\StringTok{"Humedad excesiva: revisar drenaje"}\NormalTok{)}
\end{Highlighting}
\end{Shaded}

\begin{tcolorbox}[enhanced jigsaw, arc=.35mm, bottomrule=.15mm, bottomtitle=1mm, breakable, colback=white, colbacktitle=quarto-callout-warning-color!10!white, colframe=quarto-callout-warning-color-frame, coltitle=black, left=2mm, leftrule=.75mm, opacityback=0, opacitybacktitle=0.6, rightrule=.15mm, title=\textcolor{quarto-callout-warning-color}{\faExclamationTriangle}\hspace{0.5em}{Advertencia}, titlerule=0mm, toprule=.15mm, toptitle=1mm]

\textbf{La indentación manda.} Fíjate que las líneas después de
\texttt{if} van desplazadas 4 espacios a la derecha. Ese sangrado le
dice a Python qué pertenece a cada bloque. Si lo quitas, da error. Es la
regla más famosa de Python: la sangría no es decoración, es sintaxis.

\end{tcolorbox}

\subsection{Operadores de comparación y
lógicos}\label{operadores-de-comparaciuxf3n-y-luxf3gicos}

Las condiciones se construyen con estos operadores:

\begin{longtable}[]{@{}lclcc@{}}
\caption{Operadores de comparación y
lógicos}\label{tbl-comparacion}\tabularnewline
\toprule\noalign{}
Tipo & Operador & Significado & Ejemplo & Resultado \\
\midrule\noalign{}
\endfirsthead
\toprule\noalign{}
Tipo & Operador & Significado & Ejemplo & Resultado \\
\midrule\noalign{}
\endhead
\bottomrule\noalign{}
\endlastfoot
Comparación & \texttt{==} & Igual a & \texttt{5\ ==\ 5} &
\texttt{True} \\
Comparación & \texttt{!=} & Diferente de & \texttt{5\ !=\ 3} &
\texttt{True} \\
Comparación & \texttt{\textless{}} & Menor que &
\texttt{25\ \textless{}\ 30} & \texttt{True} \\
Comparación & \texttt{\textgreater{}} & Mayor que &
\texttt{7\ \textgreater{}\ 4} & \texttt{True} \\
Comparación & \texttt{\textless{}=} & Menor o igual &
\texttt{3\ \textless{}=\ 2} & \texttt{False} \\
Comparación & \texttt{\textgreater{}=} & Mayor o igual &
\texttt{30\ \textgreater{}=\ 30} & \texttt{True} \\
Lógico & \texttt{and} & Y (ambas verdaderas) & \texttt{True\ and\ False}
& \texttt{False} \\
Lógico & \texttt{or} & O (al menos una) & \texttt{True\ or\ False} &
\texttt{True} \\
Lógico & \texttt{not} & Negación & \texttt{not\ True} &
\texttt{False} \\
\end{longtable}

\begin{tcolorbox}[enhanced jigsaw, arc=.35mm, bottomrule=.15mm, bottomtitle=1mm, breakable, colback=white, colbacktitle=quarto-callout-important-color!10!white, colframe=quarto-callout-important-color-frame, coltitle=black, left=2mm, leftrule=.75mm, opacityback=0, opacitybacktitle=0.6, rightrule=.15mm, title=\textcolor{quarto-callout-important-color}{\faExclamation}\hspace{0.5em}{Importante}, titlerule=0mm, toprule=.15mm, toptitle=1mm]

\textbf{Doble igual para comparar.} Es el error número 1 de los
principiantes: \texttt{=} asigna, \texttt{==} compara.
\texttt{if\ humedad\ =\ 30:} es error; \texttt{if\ humedad\ ==\ 30:} es
correcto.

\end{tcolorbox}

Con \texttt{and} y \texttt{or} puedes combinar condiciones, como una
alarma de invernadero que se activa por temperatura \textbf{o} por
humedad:

\begin{Shaded}
\begin{Highlighting}[]
\NormalTok{temperatura }\OperatorTok{=} \DecValTok{38}
\NormalTok{humedad }\OperatorTok{=} \DecValTok{45}

\ControlFlowTok{if}\NormalTok{ temperatura }\OperatorTok{\textgreater{}} \DecValTok{35} \KeywordTok{or}\NormalTok{ humedad }\OperatorTok{\textgreater{}} \DecValTok{80}\NormalTok{:}
    \BuiltInTok{print}\NormalTok{(}\StringTok{"ALARMA: condiciones fuera de rango"}\NormalTok{)}
\ControlFlowTok{else}\NormalTok{:}
    \BuiltInTok{print}\NormalTok{(}\StringTok{"Condiciones normales"}\NormalTok{)}
\end{Highlighting}
\end{Shaded}

\section{\texorpdfstring{Repeticiones: el ciclo
\texttt{for}}{Repeticiones: el ciclo for}}\label{repeticiones-el-ciclo-for}

Un sensor toma una lectura cada hora. Un administrador calcula la nómina
de 50 empleados. Un ingeniero evalúa 12 parcelas. En todos los casos:
\textbf{la misma operación, muchas veces}. Para eso está el ciclo
\texttt{for}.

La forma más simple usa \texttt{range()}, que genera una secuencia de
números:

\begin{Shaded}
\begin{Highlighting}[]
\ControlFlowTok{for}\NormalTok{ hora }\KeywordTok{in} \BuiltInTok{range}\NormalTok{(}\DecValTok{24}\NormalTok{):}
    \BuiltInTok{print}\NormalTok{(}\SpecialStringTok{f"Tomando lectura de la hora }\SpecialCharTok{\{}\NormalTok{hora}\SpecialCharTok{\}}\SpecialStringTok{"}\NormalTok{)}
\end{Highlighting}
\end{Shaded}

Esto ejecuta el \texttt{print} 24 veces, con \texttt{hora} tomando los
valores 0, 1, 2, \ldots, 23.

Ejemplo agro: calcular la producción total recorriendo las parcelas de
una finca:

\begin{Shaded}
\begin{Highlighting}[]
\NormalTok{producciones }\OperatorTok{=}\NormalTok{ [}\FloatTok{4.2}\NormalTok{, }\FloatTok{3.8}\NormalTok{, }\FloatTok{5.1}\NormalTok{, }\FloatTok{4.6}\NormalTok{, }\FloatTok{3.9}\NormalTok{]  }\CommentTok{\# toneladas por parcela}
\NormalTok{total }\OperatorTok{=} \DecValTok{0}

\ControlFlowTok{for}\NormalTok{ p }\KeywordTok{in}\NormalTok{ producciones:}
\NormalTok{    total }\OperatorTok{=}\NormalTok{ total }\OperatorTok{+}\NormalTok{ p}

\BuiltInTok{print}\NormalTok{(}\SpecialStringTok{f"Producción total: }\SpecialCharTok{\{}\NormalTok{total}\SpecialCharTok{\}}\SpecialStringTok{ toneladas"}\NormalTok{)}
\end{Highlighting}
\end{Shaded}

El patrón \texttt{total\ =\ total\ +\ p} se llama \textbf{acumulador}:
una variable que empieza en cero y va sumando. Lo usarás constantemente.

Ejemplo de administración: proyectar las ventas de los próximos 5 años
con crecimiento del 8 \% anual:

\begin{Shaded}
\begin{Highlighting}[]
\NormalTok{ventas }\OperatorTok{=} \DecValTok{50000000}  \CommentTok{\# ventas del año actual}

\ControlFlowTok{for}\NormalTok{ año }\KeywordTok{in} \BuiltInTok{range}\NormalTok{(}\DecValTok{1}\NormalTok{, }\DecValTok{6}\NormalTok{):}
\NormalTok{    ventas }\OperatorTok{=}\NormalTok{ ventas }\OperatorTok{*} \FloatTok{1.08}
    \BuiltInTok{print}\NormalTok{(}\SpecialStringTok{f"Año }\SpecialCharTok{\{}\NormalTok{a}\SpecialCharTok{ñ}\NormalTok{o}\SpecialCharTok{\}}\SpecialStringTok{: $}\SpecialCharTok{\{}\NormalTok{ventas}\SpecialCharTok{:,.0f\}}\SpecialStringTok{"}\NormalTok{)}
\end{Highlighting}
\end{Shaded}

Fíjate en el truco \texttt{\{ventas:,.0f\}} dentro de la f-string:
formatea el número con separador de miles y sin decimales.

\section{\texorpdfstring{Repeticiones con condición: el ciclo
\texttt{while}}{Repeticiones con condición: el ciclo while}}\label{repeticiones-con-condiciuxf3n-el-ciclo-while}

El \texttt{for} repite un número conocido de veces. El \texttt{while}
repite \textbf{mientras una condición sea verdadera} --- cuando no sabes
cuántas vueltas harán falta:

\begin{Shaded}
\begin{Highlighting}[]
\NormalTok{tanque }\OperatorTok{=} \DecValTok{0}      \CommentTok{\# litros}
\NormalTok{capacidad }\OperatorTok{=} \DecValTok{1000}

\ControlFlowTok{while}\NormalTok{ tanque }\OperatorTok{\textless{}}\NormalTok{ capacidad:}
\NormalTok{    tanque }\OperatorTok{=}\NormalTok{ tanque }\OperatorTok{+} \DecValTok{150}   \CommentTok{\# la bomba llena 150 litros por ciclo}
    \BuiltInTok{print}\NormalTok{(}\SpecialStringTok{f"Nivel actual: }\SpecialCharTok{\{}\NormalTok{tanque}\SpecialCharTok{\}}\SpecialStringTok{ litros"}\NormalTok{)}

\BuiltInTok{print}\NormalTok{(}\StringTok{"Tanque lleno"}\NormalTok{)}
\end{Highlighting}
\end{Shaded}

\begin{tcolorbox}[enhanced jigsaw, arc=.35mm, bottomrule=.15mm, bottomtitle=1mm, breakable, colback=white, colbacktitle=quarto-callout-warning-color!10!white, colframe=quarto-callout-warning-color-frame, coltitle=black, left=2mm, leftrule=.75mm, opacityback=0, opacitybacktitle=0.6, rightrule=.15mm, title=\textcolor{quarto-callout-warning-color}{\faExclamationTriangle}\hspace{0.5em}{Advertencia}, titlerule=0mm, toprule=.15mm, toptitle=1mm]

Si la condición de un \texttt{while} jamás se hace falsa, el ciclo corre
para siempre (ciclo infinito). En Colab lo detienes con el botón ⏹️.
Asegúrate siempre de que algo dentro del ciclo acerque la condición a su
fin.

\end{tcolorbox}

\section{Empaquetar: las funciones}\label{empaquetar-las-funciones}

Ya usaste funciones que vienen con Python: \texttt{print()},
\texttt{input()}, \texttt{int()}. Ahora crearás las tuyas. Una función
es un procedimiento con nombre: le das entradas (\textbf{parámetros}),
hace su trabajo y te devuelve un resultado (\textbf{retorno}).

\begin{Shaded}
\begin{Highlighting}[]
\KeywordTok{def}\NormalTok{ valor\_futuro(capital, tasa, años):}
\NormalTok{    vf }\OperatorTok{=}\NormalTok{ capital }\OperatorTok{*}\NormalTok{ (}\DecValTok{1} \OperatorTok{+}\NormalTok{ tasa) }\OperatorTok{**}\NormalTok{ años}
    \ControlFlowTok{return}\NormalTok{ vf}
\end{Highlighting}
\end{Shaded}

\begin{itemize}
\tightlist
\item
  \texttt{def} le dice a Python ``voy a definir una función''.
\item
  Entre paréntesis van los parámetros: lo que la función necesita
  recibir.
\item
  \texttt{return} entrega el resultado y termina la función.
\end{itemize}

Se usa así, todas las veces que quieras:

\begin{Shaded}
\begin{Highlighting}[]
\NormalTok{inversion1 }\OperatorTok{=}\NormalTok{ valor\_futuro(}\DecValTok{1000000}\NormalTok{, }\FloatTok{0.12}\NormalTok{, }\DecValTok{5}\NormalTok{)}
\NormalTok{inversion2 }\OperatorTok{=}\NormalTok{ valor\_futuro(}\DecValTok{5000000}\NormalTok{, }\FloatTok{0.08}\NormalTok{, }\DecValTok{10}\NormalTok{)}

\BuiltInTok{print}\NormalTok{(}\SpecialStringTok{f"Proyecto A: $}\SpecialCharTok{\{}\NormalTok{inversion1}\SpecialCharTok{:,.0f\}}\SpecialStringTok{"}\NormalTok{)}
\BuiltInTok{print}\NormalTok{(}\SpecialStringTok{f"Proyecto B: $}\SpecialCharTok{\{}\NormalTok{inversion2}\SpecialCharTok{:,.0f\}}\SpecialStringTok{"}\NormalTok{)}
\end{Highlighting}
\end{Shaded}

La ventaja enorme: la fórmula se escribe \textbf{una sola vez}. Si
mañana hay que corregirla, se corrige en un solo lugar.

Otro ejemplo, ahora de mecatrónica: una función que clasifica la lectura
de un sensor de temperatura:

\begin{Shaded}
\begin{Highlighting}[]
\KeywordTok{def}\NormalTok{ clasificar\_temperatura(t):}
    \ControlFlowTok{if}\NormalTok{ t }\OperatorTok{\textless{}} \DecValTok{15}\NormalTok{:}
        \ControlFlowTok{return} \StringTok{"frío"}
    \ControlFlowTok{elif}\NormalTok{ t }\OperatorTok{\textless{}=} \DecValTok{30}\NormalTok{:}
        \ControlFlowTok{return} \StringTok{"óptimo"}
    \ControlFlowTok{else}\NormalTok{:}
        \ControlFlowTok{return} \StringTok{"caliente"}

\BuiltInTok{print}\NormalTok{(clasificar\_temperatura(}\DecValTok{22}\NormalTok{))}
\BuiltInTok{print}\NormalTok{(clasificar\_temperatura(}\DecValTok{38}\NormalTok{))}
\end{Highlighting}
\end{Shaded}

\section{Ejercicios}\label{ejercicios-2}

\begin{enumerate}
\def\labelenumi{\arabic{enumi}.}
\tightlist
\item
  Escribe un \texttt{if} que imprima ``Crédito aprobado'' si el puntaje
  crediticio es mayor o igual a 700, y ``Crédito rechazado'' en caso
  contrario. Prueba con 650 y con 720.
\item
  Un motor debe apagarse si la temperatura supera 90 °C \textbf{y} el
  voltaje es inestable (variable booleana \texttt{voltaje\_inestable}).
  Escribe la condición con \texttt{and} y pruébala con distintos
  valores.
\item
  Con un \texttt{for} y un acumulador, suma los costos de una lista de
  insumos: \texttt{{[}120000,\ 85000,\ 210000,\ 45000{]}}. Imprime el
  total.
\item
  Escribe una función \texttt{area\_parcela(largo,\ ancho)} que devuelva
  el área en metros cuadrados, y úsala para tres parcelas distintas.
\item
  Escribe una función \texttt{es\_rentable(ingresos,\ costos)} que
  devuelva \texttt{True} si los ingresos superan a los costos y
  \texttt{False} si no.
\item
  \textbf{Mini-reto (agro + economía):} tienes las producciones de 6
  parcelas: \texttt{{[}4.2,\ 2.8,\ 5.1,\ 3.3,\ 4.6,\ 2.1{]}} toneladas.
  Con un \texttt{for} y un contador, calcula cuántas parcelas produjeron
  menos de 3 toneladas (bajo rendimiento) e imprime el porcentaje que
  representan del total. Luego empaqueta tu solución en una función
  \texttt{porcentaje\_bajo\_rendimiento(producciones,\ umbral)}.
\end{enumerate}

\section{Resumen}\label{resumen-2}

\begin{itemize}
\tightlist
\item
  \texttt{if} / \texttt{elif} / \texttt{else} ejecutan bloques según
  condiciones; la \textbf{indentación define los bloques}.
\item
  Comparar con \texttt{==}, asignar con \texttt{=}. Combinar condiciones
  con \texttt{and}, \texttt{or}, \texttt{not}.
\item
  \texttt{for} repite un número conocido de veces; \texttt{while} repite
  mientras la condición sea verdadera.
\item
  El \textbf{acumulador} (\texttt{total\ =\ total\ +\ x}) es el patrón
  básico para sumar en ciclos.
\item
  Las \textbf{funciones} (\texttt{def} \ldots{} \texttt{return})
  empaquetan procedimientos para reutilizarlos sin duplicar código.
\end{itemize}

\bookmarksetup{startatroot}

\chapter{Estructuras de datos: listas, tuplas y
diccionarios}\label{estructuras-de-datos-listas-tuplas-y-diccionarios}

Una variable guarda \textbf{un} valor. Pero en la práctica manejas
colecciones: las lecturas de un sensor durante el día, los precios de 20
productos, las coordenadas de las parcelas. Para eso existen las
estructuras de datos. En este capítulo verás las tres esenciales:
\textbf{listas}, \textbf{tuplas} y \textbf{diccionarios}.

\section{Listas: colecciones ordenadas que puedes
modificar}\label{listas-colecciones-ordenadas-que-puedes-modificar}

Una lista es una secuencia de valores entre corchetes \texttt{{[}{]}},
separados por comas:

\begin{Shaded}
\begin{Highlighting}[]
\NormalTok{lecturas }\OperatorTok{=}\NormalTok{ [}\FloatTok{22.5}\NormalTok{, }\FloatTok{23.1}\NormalTok{, }\FloatTok{24.0}\NormalTok{, }\FloatTok{23.8}\NormalTok{, }\FloatTok{25.2}\NormalTok{]  }\CommentTok{\# temperaturas de un sensor}
\end{Highlighting}
\end{Shaded}

\subsection{Acceder a un elemento: los índices empiezan en
cero}\label{acceder-a-un-elemento-los-uxedndices-empiezan-en-cero}

Cada elemento tiene una posición llamada \textbf{índice}. La regla que
confunde a todo principiante: \textbf{el primer elemento es el índice
0}, no el 1:

\begin{Shaded}
\begin{Highlighting}[]
\NormalTok{lecturas }\OperatorTok{=}\NormalTok{ [}\FloatTok{22.5}\NormalTok{, }\FloatTok{23.1}\NormalTok{, }\FloatTok{24.0}\NormalTok{, }\FloatTok{23.8}\NormalTok{, }\FloatTok{25.2}\NormalTok{]}

\BuiltInTok{print}\NormalTok{(lecturas[}\DecValTok{0}\NormalTok{])   }\CommentTok{\# primera lectura}
\BuiltInTok{print}\NormalTok{(lecturas[}\DecValTok{2}\NormalTok{])   }\CommentTok{\# tercera lectura}
\BuiltInTok{print}\NormalTok{(lecturas[}\OperatorTok{{-}}\DecValTok{1}\NormalTok{])  }\CommentTok{\# última lectura (el {-}1 cuenta desde el final)}
\end{Highlighting}
\end{Shaded}

\subsection{Modificar, agregar y
contar}\label{modificar-agregar-y-contar}

Las listas son \textbf{mutables}: puedes cambiarlas después de crearlas:

\begin{Shaded}
\begin{Highlighting}[]
\NormalTok{precios }\OperatorTok{=}\NormalTok{ [}\DecValTok{9500}\NormalTok{, }\DecValTok{8200}\NormalTok{, }\DecValTok{10500}\NormalTok{]}

\NormalTok{precios[}\DecValTok{0}\NormalTok{] }\OperatorTok{=} \DecValTok{9800}        \CommentTok{\# modificar: el primer precio cambió}
\NormalTok{precios.append(}\DecValTok{7600}\NormalTok{)     }\CommentTok{\# agregar al final}
\BuiltInTok{print}\NormalTok{(precios)}
\BuiltInTok{print}\NormalTok{(}\BuiltInTok{len}\NormalTok{(precios))      }\CommentTok{\# cuántos elementos tiene}
\end{Highlighting}
\end{Shaded}

\subsection{Operaciones útiles sin escribir
ciclos}\label{operaciones-uxfatiles-sin-escribir-ciclos}

Python trae funciones que resumen una lista en una línea:

\begin{Shaded}
\begin{Highlighting}[]
\NormalTok{producciones }\OperatorTok{=}\NormalTok{ [}\FloatTok{4.2}\NormalTok{, }\FloatTok{3.8}\NormalTok{, }\FloatTok{5.1}\NormalTok{, }\FloatTok{4.6}\NormalTok{, }\FloatTok{3.9}\NormalTok{]}

\BuiltInTok{print}\NormalTok{(}\BuiltInTok{sum}\NormalTok{(producciones))   }\CommentTok{\# total}
\BuiltInTok{print}\NormalTok{(}\BuiltInTok{max}\NormalTok{(producciones))   }\CommentTok{\# la mayor}
\BuiltInTok{print}\NormalTok{(}\BuiltInTok{min}\NormalTok{(producciones))   }\CommentTok{\# la menor}
\BuiltInTok{print}\NormalTok{(}\BuiltInTok{len}\NormalTok{(producciones))   }\CommentTok{\# cuántas parcelas}
\BuiltInTok{print}\NormalTok{(}\BuiltInTok{sum}\NormalTok{(producciones) }\OperatorTok{/} \BuiltInTok{len}\NormalTok{(producciones))  }\CommentTok{\# promedio}
\end{Highlighting}
\end{Shaded}

\subsection{Rebanadas (slicing): tomar una
parte}\label{rebanadas-slicing-tomar-una-parte}

Con \texttt{{[}inicio:fin{]}} tomas un pedazo de la lista (incluye el
inicio, \textbf{no} incluye el fin):

\begin{Shaded}
\begin{Highlighting}[]
\NormalTok{lecturas }\OperatorTok{=}\NormalTok{ [}\FloatTok{22.5}\NormalTok{, }\FloatTok{23.1}\NormalTok{, }\FloatTok{24.0}\NormalTok{, }\FloatTok{23.8}\NormalTok{, }\FloatTok{25.2}\NormalTok{]}

\BuiltInTok{print}\NormalTok{(lecturas[}\DecValTok{1}\NormalTok{:}\DecValTok{3}\NormalTok{])   }\CommentTok{\# del índice 1 al 2}
\BuiltInTok{print}\NormalTok{(lecturas[:}\DecValTok{2}\NormalTok{])    }\CommentTok{\# las dos primeras}
\BuiltInTok{print}\NormalTok{(lecturas[}\DecValTok{3}\NormalTok{:])    }\CommentTok{\# del índice 3 al final}
\end{Highlighting}
\end{Shaded}

\subsection{Recorrer una lista}\label{recorrer-una-lista}

Ya lo hiciste en el capítulo anterior: el \texttt{for} nació para esto:

\begin{Shaded}
\begin{Highlighting}[]
\NormalTok{ventas }\OperatorTok{=}\NormalTok{ [}\FloatTok{12.5}\NormalTok{, }\FloatTok{15.2}\NormalTok{, }\FloatTok{9.8}\NormalTok{, }\FloatTok{18.1}\NormalTok{]  }\CommentTok{\# millones por trimestre}

\ControlFlowTok{for}\NormalTok{ v }\KeywordTok{in}\NormalTok{ ventas:}
    \ControlFlowTok{if}\NormalTok{ v }\OperatorTok{\textless{}} \DecValTok{10}\NormalTok{:}
        \BuiltInTok{print}\NormalTok{(}\SpecialStringTok{f"Trimestre bajo: }\SpecialCharTok{\{}\NormalTok{v}\SpecialCharTok{\}}\SpecialStringTok{"}\NormalTok{)}
\end{Highlighting}
\end{Shaded}

\section{Tuplas: colecciones que no deben
cambiar}\label{tuplas-colecciones-que-no-deben-cambiar}

Una tupla se escribe con paréntesis \texttt{()} y es \textbf{inmutable}:
una vez creada, no se puede modificar. Sirve para datos que
conceptualmente van juntos y no deben alterarse, como las coordenadas de
una parcela:

\begin{Shaded}
\begin{Highlighting}[]
\NormalTok{ubicacion }\OperatorTok{=}\NormalTok{ (}\FloatTok{3.42}\NormalTok{, }\OperatorTok{{-}}\FloatTok{76.53}\NormalTok{)  }\CommentTok{\# latitud, longitud}

\BuiltInTok{print}\NormalTok{(ubicacion[}\DecValTok{0}\NormalTok{])  }\CommentTok{\# latitud}
\end{Highlighting}
\end{Shaded}

Si intentas modificarla, Python te protege con un error:

\begin{Shaded}
\begin{Highlighting}[]
\NormalTok{ubicacion[}\DecValTok{0}\NormalTok{] }\OperatorTok{=} \FloatTok{4.0}   \CommentTok{\# TypeError: las tuplas no se modifican}
\end{Highlighting}
\end{Shaded}

Regla práctica: si los datos representan algo fijo (coordenadas, una
fecha, un registro que no debe tocarse), usa tupla. Si van a cambiar
(lecturas, precios, inventarios), usa lista.

\section{Diccionarios: datos con nombre, no con
posición}\label{diccionarios-datos-con-nombre-no-con-posiciuxf3n}

En una lista accedes por posición (\texttt{precios{[}0{]}}). Pero ¿y si
quieres acceder \textbf{por nombre}? Para eso están los diccionarios:
pares de \textbf{clave} y \textbf{valor} entre llaves \texttt{\{\}}:

\begin{Shaded}
\begin{Highlighting}[]
\NormalTok{precios }\OperatorTok{=}\NormalTok{ \{}
    \StringTok{"café"}\NormalTok{: }\DecValTok{9500}\NormalTok{,}
    \StringTok{"cacao"}\NormalTok{: }\DecValTok{8200}\NormalTok{,}
    \StringTok{"banano"}\NormalTok{: }\DecValTok{7600}
\NormalTok{\}}

\BuiltInTok{print}\NormalTok{(precios[}\StringTok{"café"}\NormalTok{])   }\CommentTok{\# acceder por clave}
\end{Highlighting}
\end{Shaded}

Agregar y modificar es asignar a una clave:

\begin{Shaded}
\begin{Highlighting}[]
\NormalTok{precios[}\StringTok{"banano"}\NormalTok{] }\OperatorTok{=} \DecValTok{7800}      \CommentTok{\# modificar}
\NormalTok{precios[}\StringTok{"panela"}\NormalTok{] }\OperatorTok{=} \DecValTok{5200}      \CommentTok{\# agregar un producto nuevo}
\BuiltInTok{print}\NormalTok{(precios)}
\end{Highlighting}
\end{Shaded}

Recorrer un diccionario con \texttt{.items()} te da clave y valor a la
vez:

\begin{Shaded}
\begin{Highlighting}[]
\NormalTok{inventario }\OperatorTok{=}\NormalTok{ \{}\StringTok{"motobombas"}\NormalTok{: }\DecValTok{3}\NormalTok{, }\StringTok{"sensores\_humedad"}\NormalTok{: }\DecValTok{12}\NormalTok{, }\StringTok{"valvulas"}\NormalTok{: }\DecValTok{25}\NormalTok{\}}

\ControlFlowTok{for}\NormalTok{ equipo, cantidad }\KeywordTok{in}\NormalTok{ inventario.items():}
    \ControlFlowTok{if}\NormalTok{ cantidad }\OperatorTok{\textless{}} \DecValTok{5}\NormalTok{:}
        \BuiltInTok{print}\NormalTok{(}\SpecialStringTok{f"Reponer }\SpecialCharTok{\{}\NormalTok{equipo}\SpecialCharTok{\}}\SpecialStringTok{: quedan }\SpecialCharTok{\{}\NormalTok{cantidad}\SpecialCharTok{\}}\SpecialStringTok{"}\NormalTok{)}
\end{Highlighting}
\end{Shaded}

\subsection{Un diccionario puede guardar una ficha
completa}\label{un-diccionario-puede-guardar-una-ficha-completa}

Los diccionarios brillan cuando un valor es a su vez otra estructura:

\begin{Shaded}
\begin{Highlighting}[]
\NormalTok{empleado }\OperatorTok{=}\NormalTok{ \{}
    \StringTok{"nombre"}\NormalTok{: }\StringTok{"María Pérez"}\NormalTok{,}
    \StringTok{"cargo"}\NormalTok{: }\StringTok{"Analista"}\NormalTok{,}
    \StringTok{"salario"}\NormalTok{: }\DecValTok{2800000}\NormalTok{,}
    \StringTok{"areas"}\NormalTok{: [}\StringTok{"compras"}\NormalTok{, }\StringTok{"logística"}\NormalTok{]}
\NormalTok{\}}

\BuiltInTok{print}\NormalTok{(empleado[}\StringTok{"nombre"}\NormalTok{])}
\BuiltInTok{print}\NormalTok{(empleado[}\StringTok{"salario"}\NormalTok{] }\OperatorTok{*} \DecValTok{12}\NormalTok{)   }\CommentTok{\# salario anual}
\end{Highlighting}
\end{Shaded}

Este patrón --- una ficha con campos nombrados --- es la base de casi
todo el manejo de datos que verás en el capítulo de pandas.

\section{¿Cuál estructura usar?}\label{cuuxe1l-estructura-usar}

{\def\LTcaptype{none} % do not increment counter
\begin{longtable}[]{@{}
  >{\raggedright\arraybackslash}p{(\linewidth - 8\tabcolsep) * \real{0.1967}}
  >{\centering\arraybackslash}p{(\linewidth - 8\tabcolsep) * \real{0.1639}}
  >{\centering\arraybackslash}p{(\linewidth - 8\tabcolsep) * \real{0.2459}}
  >{\raggedright\arraybackslash}p{(\linewidth - 8\tabcolsep) * \real{0.1967}}
  >{\raggedright\arraybackslash}p{(\linewidth - 8\tabcolsep) * \real{0.1967}}@{}}
\toprule\noalign{}
\begin{minipage}[b]{\linewidth}\raggedright
Estructura
\end{minipage} & \begin{minipage}[b]{\linewidth}\centering
Sintaxis
\end{minipage} & \begin{minipage}[b]{\linewidth}\centering
¿Se modifica?
\end{minipage} & \begin{minipage}[b]{\linewidth}\raggedright
Acceso por
\end{minipage} & \begin{minipage}[b]{\linewidth}\raggedright
Úsala para
\end{minipage} \\
\midrule\noalign{}
\endhead
\bottomrule\noalign{}
\endlastfoot
Lista & \texttt{{[}1,\ 2,\ 3{]}} & Sí & Posición & Series de mediciones,
colecciones que crecen \\
Tupla & \texttt{(1,\ 2)} & No & Posición & Datos fijos que van juntos
(coordenadas) \\
Diccionario & \texttt{\{"a":\ 1\}} & Sí & Nombre (clave) & Fichas,
inventarios, catálogos, configuración \\
\end{longtable}
}

\section{Ejercicios}\label{ejercicios-3}

\begin{enumerate}
\def\labelenumi{\arabic{enumi}.}
\tightlist
\item
  Crea una lista con las temperaturas de 7 días. Imprime la máxima, la
  mínima y el promedio (con \texttt{sum} y \texttt{len}).
\item
  De la lista
  \texttt{costos\ =\ {[}45000,\ 120000,\ 38000,\ 90000,\ 67000{]}},
  imprime solo los costos mayores a 50000 usando un \texttt{for} y un
  \texttt{if}.
\item
  Crea una tupla con la coordenada de tu ciudad y otra con la de una
  finca. Imprime la distancia en latitud (diferencia simple del primer
  elemento).
\item
  Crea un diccionario con 4 productos agrícolas y sus precios. Sube
  todos los precios un 10 \% recorriéndolo con \texttt{.items()} y
  reasignando cada valor.
\item
  Crea un diccionario que represente un sensor con las claves
  \texttt{id}, \texttt{tipo}, \texttt{parcela} y \texttt{lecturas} (una
  lista de 5 valores). Imprime el promedio de sus lecturas accediendo a
  través del diccionario.
\item
  \textbf{Mini-reto (administración):} tienes esta lista de ventas
  mensuales (millones):
  \texttt{{[}42,\ 38,\ 55,\ 61,\ 47,\ 39,\ 52,\ 58,\ 63,\ 49,\ 44,\ 57{]}}.
  Sin usar funciones hechas, solo un \texttt{for} y acumuladores:
  calcula el promedio anual, cuenta cuántos meses superaron el promedio
  e imprime ambos resultados. Luego, guarda en un diccionario
  \texttt{resumen} las claves \texttt{"promedio"} y
  \texttt{"meses\_sobre\_promedio"}.
\end{enumerate}

\section{Resumen}\label{resumen-3}

\begin{itemize}
\tightlist
\item
  \textbf{Lista} \texttt{{[}{]}}: ordenada y modificable; el caballo de
  batalla. Índices desde \textbf{0}.
\item
  \texttt{append()} agrega, \texttt{len()} cuenta,
  \texttt{sum()/max()/min()} resumen, \texttt{{[}i:j{]}} rebana.
\item
  \textbf{Tupla} \texttt{()}: igual a la lista pero inmutable; para
  datos fijos.
\item
  \textbf{Diccionario} \texttt{\{\}}: pares clave-valor; acceso por
  nombre con \texttt{d{[}"clave"{]}}, recorrido con \texttt{.items()}.
\item
  Elegir bien la estructura es la mitad del programa.
\end{itemize}

\bookmarksetup{startatroot}

\chapter{Capítulo 5: NumPy y
gráficas}\label{capuxedtulo-5-numpy-y-gruxe1ficas}

En construcción.

\bookmarksetup{startatroot}

\chapter{Capítulo 6: pandas}\label{capuxedtulo-6-pandas}

En construcción.

\bookmarksetup{startatroot}

\chapter{Capítulo 7: Economía y
administración}\label{capuxedtulo-7-economuxeda-y-administraciuxf3n}

En construcción.

\bookmarksetup{startatroot}

\chapter{Capítulo 8: Introducción a big
data}\label{capuxedtulo-8-introducciuxf3n-a-big-data}

En construcción.




\end{document}
